
\subsubsection{5.1 Main Result: SE\_N5 and min\_logprob Are Near-Orthogonal (RQ1)}

\textbf{Pearson |r|(SE\_N5, min\_logprob) = 0.026} at N=2500 on TriviaQA dev with Llama-3.1-8B-Instruct. This value is far below both the independence threshold (0.70) and the reparameterization boundary (0.85), meeting Gate 1 with substantial margin.

Figure 1 (\texttt{h-e1/figures/scatter\textbackslash \_se\textbackslash \_vs\textbackslash \_minlogprob.png}) visualizes this independence: 2500 data points show no discernible linear or monotone pattern between SE and min\_logprob values, with correctness labels scattered throughout the joint space. The ABANDON threshold was not approached, ruling out SE as a monotone reparameterization of log-probability.

Figure 2 (\texttt{h-e1/figures/correlation\textbackslash \_heatmap.png}) shows the full Pearson correlation matrix between SE\_N5, min\_logprob, response length (L), and LM-judge correctness. The near-zero cross-correlation between SE and min\_logprob (|r|=0.026) stands in contrast to the moderate negative correlations that both signals exhibit individually with correctness --- each carries information about correctness, but through independent mechanisms.

\begin{table*}[htbp]
\centering
\resizebox{\dimexpr 2\columnwidth+\columnsep\relax}{!}{%
\begin{tabular}{llllll}
\toprule
Metric & Value & Threshold & Status &  &  \\
\midrule
Pearson \ & r\ & (SE\_N5, min\_logprob) & **0.026** & < 0.70 & **PASS** \\
ABANDON threshold & --- & > 0.85 & NOT triggered &  &  \\
SE variance & 0.100 & > 0.01 & PASS (non-degenerate) &  &  \\
Fraction degenerate & 0.000 & = 0 & PASS &  &  \\
min\_logprob mean & $-2.435$ & < 0 & PASS &  &  \\
\bottomrule
\end{tabular}
}
\end{table*}

The mechanism reality checks confirm both signals are well-behaved: SE variance = 0.100 (well above non-degeneracy threshold), fraction\_degenerate = 0 (no prompts where all 5 stochastic samples collapse to the same semantic cluster), and min\_logprob mean = $-2.435$ (all values are valid negative log-probabilities).

\subsubsection{5.2 Conditional Independence Analysis (RQ2)}

Figure 3 (\texttt{h-e1/figures/gate\textbackslash \_metrics.png}) summarizes the gate evaluation: the Pearson |r|=0.026 bar clears the 0.70 threshold with wide margin, while the partial $R^{2}=0.0005$ bar falls below the 0.02 threshold (LRT $chi^{2}=1.632$, df=2, p=0.442, non-significant). Figure 4 (\texttt{h-e1/figures/lr\textbackslash \_coefficients.png}) shows the conditional logistic regression coefficients with 95\% CI. The SE coefficient is positive (correct direction: higher SE predicts lower correctness), but the confidence interval crosses zero at N=2500 --- consistent with a powered test yielding a genuine null result.

\begin{table}[htbp]
\centering
\resizebox{\columnwidth}{!}{%
\begin{tabular}{llll}
\toprule
LR Term & Direction & Significance & p-value \\
\midrule
min\_logprob & positive & significant (CI above zero) & < 0.05 \\
SE\_N5 & positive & not significant (CI crosses zero) & > 0.05 \\
Partial $R^{2}(SE)$ & --- & **0.0005** (target $\geq$ 0.02) & p = 0.442 \\
\bottomrule
\end{tabular}
}
\end{table}

The partial $R^2$ = 0.0005 at N=2500 does not meet the pre-registered $\geq$ 0.02 threshold. The LRT provides no evidence that SE adds conditional predictive power beyond min\_logprob in a linear setting.

\textbf{Gate outcome:} The pre-registered pipeline gate classified this result as \texttt{EXPLORE\textbackslash \_N10} (gate\_pass: false, reason: ''abs(r)=0.026 or partial\_r2=0.0005 marginal --- retry N=10''). The gate did not pass. The EXPLORE\_N10 decision indicates the pipeline recommends extending to N=10 stochastic samples per prompt. A standard power analysis for logistic regression at N=2500 with binary outcome prevalence of 34.8\% indicates that N=2500 substantially exceeds the sample size required to detect partial $R^{2}\geq 0.02$ at $\alpha =0.05$ with >0.80 power (estimated power >0.99 at the pre-registered effect size). The gate's EXPLORE\_N10 flag reflects the pipeline's conservative marginal-uncertainty criterion applied to both metrics jointly, not a determination that the statistical test was underpowered for the pre-registered threshold. These two characterizations are not contradictory: the gate reflects pipeline confidence policy; the power characterization reflects adequacy of N=2500 to detect the pre-registered effect size.

\subsubsection{5.3 Circularity Control Validation (RQ3)}

\textbf{Spearman $\rho$(SE, LM-judge correctness) = $-0.026$}, well within the |$\rho$| < 0.40 circularity threshold. The negative sign is directionally consistent: higher SE (more semantic uncertainty) correlates weakly with lower correctness. The small magnitude confirms that the cross-model judge design prevents circular agreement between the NLI-based SE signal and the judge.

Note: the Spearman |$\rho$|=0.026 and the Pearson |r|=0.026 reported in Section 5.1 share the same rounded magnitude; these are numerically distinct statistics computed on different variable pairs (SE vs. min\_logprob for Pearson r; SE vs. judge labels for Spearman $\rho$), using algebraically different correlation measures. The coincidence is real --- both reflect a near-zero correlation structure in this dataset --- and is not a copy-paste artifact.

The correctness rate of 34.8\% provides meaningful error headroom for UQ discrimination.

\subsubsection{5.4 Analysis: Why |r|=0.026 Is Substantially Lower Than Prior Work}

Gabriel (2026) reports r = 0.54--0.76 between \textbf{first-token} confidence and semantic agreement on Llama-3.1-8B, TriviaQA. The min\_logprob result (|r|=0.026) is approximately $20\times$ lower. This contrast is mechanistically explained: min\_logprob is the minimum over the full greedy decode sequence (capturing late-position tokens where factual information is concentrated), not the first token. Late-position tokens can be highly uncertain even when the first token is confident, breaking the correlation that exists for first-token signals.

\subsubsection{5.5 Summary}

\begin{table*}[htbp]
\centering
\resizebox{\dimexpr 2\columnwidth+\columnsep\relax}{!}{%
\begin{tabular}{llllll}
\toprule
RQ & Question & Finding & Status &  &  \\
\midrule
RQ1 & Are SE and min\_logprob independent? & Pearson \ & r\ & =0.026 --- near-orthogonal & **CONFIRMED** \\
RQ2 & Does SE add conditional predictive power? & Partial $R^{2}=0.0005$, LRT p=0.442 at N=2500 --- null result (gate=EXPLORE\_N10) & **NOT CONFIRMED** &  &  \\
RQ3 & Is cross-model judge low-circularity? & Spearman $\rho=-0.026$ & **CONFIRMED** &  &  \\
\bottomrule
\end{tabular}
}
\end{table*}

